\subsection{Hypercharge as a W33 duad vector and the five-dimensional neutral center cokernel}
\label{sec:20260920-hypercharge-duad-center-cokernel}

The flagship global Standard-Model-side root algebra is
\[
A_2+A_1+\mathfrak u(1)^5,
\]
and the orthogonal rank-five lattice inside \(E_8\) is exactly \(A_5\).
With standard \(A_5\) coordinates
\[
Q(A_5)=\{x\in\mathbb Z^6:\sum_i x_i=0\},
\]
the canonical hypercharge vector is
\[
Y=-\omega_1=\frac16(-5,1,1,1,1,1),
\qquad
Y^2=\frac56.
\]

The existing \(S_6\)-duad gauge module admits the normalized intertwiner
\[
\Psi(x)_{\{i,j\}}=\frac{x_i+x_j}{2}.
\]
Unlike the earlier unnormalized map, \(\Psi\) is an isometry: the five simple
roots \(e_i-e_{i+1}\) have exactly the \(A_5\) Cartan Gram matrix in the
fifteen-duad carrier.  Hypercharge therefore has the explicit W33-coordinate
form
\[
\Psi(Y)_{\{0,i\}}=-\frac13\quad(i=1,\ldots,5),
\qquad
\Psi(Y)_{\{i,j\}}=\frac16\quad(1\le i<j\le5).
\]
Thus the discriminant generator is a five-versus-ten duad pattern.  Every
syntheme sum vanishes, \(\|\Psi(Y)\|^2=5/6\), and
\[
6\Psi(Y)=(-2)^5(+1)^{10}
\]
in the star/nonstar ordering.  Its stabilizer in \(S_6\) is exactly the
\(S_5=W(A_4)\) fixing the distinguished vertex.

The hypercharge-orthogonal integral sublattice is
\[
A_4=\{x_0=0,\ \sum_{i=1}^{5}x_i=0\}.
\]
Writing
\[
y_i=\Psi(x)_{\{0,i\}}=\frac{x_i}{2},
\]
this four-dimensional sector is simply
\[
\sum_{i=1}^{5}y_i=0,
\]
and all ten nonstar coordinates are reconstructed by
\[
\Psi(x)_{\{i,j\}}=y_i+y_j.
\]
Hence the four extra Abelian directions in the flagship have an exact W33
realization as zero-sum fluctuations on the five-edge duad star.

Independently, the physical \(E_8\) holonomy bracket tensor gives the neutral
sector
\[
\dim\mathfrak g_{000}=16,
\qquad
\operatorname{rank}[\mathfrak g_{000},\mathfrak g_{000}]=11.
\]
Its five-dimensional cokernel is therefore precisely the center
\(\mathfrak u(1)^5\).  The lattice result refines this as
\[
16=11+1+4,
\]
namely
\[
\mathfrak{su}(3)\oplus\mathfrak{su}(2)
\quad+\quad
\mathfrak u(1)_Y
\quad+\quad
\mathfrak u(1)^4_{A_4}.
\]
A hypercharge-preserving singlet vacuum removes all extra Abelian factors only
if its charge projections span rank four on this \(A_4\) sector.

There is an important cardinality firewall.  The five duads in the
hypercharge star have natural stabilizer \(S_5\), whereas the five-state
\(CZ_3\) zero-phase cross has classical frame stabilizer \(D_8\).  Their equal
size does not define a canonical intertwiner.

The executable certificate is
\texttt{data/w33\_hypercharge\_duad\_center\_cokernel.json}.
